\scanpage{29}
% archive sheet number: 28
élément de $K$, on voit que le choix ainsi fait permet d'écrire
\[
 \alpha_\infty(u)=\frac{1}{2\sqrt{-1}\pi}v
\]
où $v$ est le générateur canonique de $H^1(\mathbb P^1,\Omega^1)$.

Pour les choix \unsure{différents} de topologies sur $K$ faisant de $K$ un corps valué complet \unsure{arch.}, on trouve ainsi tous les éléments possibles de $H^1(\unsure{P|\mathbb P}^1,\Omega^1)$.

Si $K$ est un corps \deleted{top.} isom à $\mathbb R$ [l'isom étant d'ailleurs \uline{unique}], alors on définit encore, par choix de $\mathbb C$ comme \deleted{\unsure{corps top.|compactification}} clôture algébrique \unsure{de|à} $\mathbb R$, un foncteur $\unsure{\xi|\zeta}=\xi_K$ à valeurs en les modules sur $\mathbb Z$, \deleted{et \unsure{un|une} \unsure{des} \unsure{isom|isom.}} \inserted{mais avec op. de $\mathbb Z/2\mathbb Z=\operatorname{Gal}(\overline{\mathbb R}/\mathbb R)$} Le sous $\deleted{\unsure{\xi}}\mathbb Z$-réseau \inserted{$\xi(M)$} que nous obtenons dans $\prod\limits_{\unsure{\ell|l}} M_\ell\times M_\infty{}_{\mathbb C}$ \inserted{\unsure{relatif|valeurs} : $\mathbb R\hookrightarrow\mathbb C$} est alors stable par ce groupe $\mathbb Z/2\mathbb Z$

